% chapters/16/praktikum.tex
\begin{praktikumbox}
Praktikum ini adalah bagian pengujian dari \textbf{Proyek Akhir} pada RPS:
menyusun \istilah{checklist} fungsional, menjalankan pemeriksaan otomatis,
melaksanakan \istilah{playtest} antar kelompok, mencatat temuan, memperbaiki,
lalu menyajikan demo lima sampai sepuluh menit.

\begin{langkah}
  \item \textbf{Tuliskan \istilah{checklist} fungsionalnya.} Bacalah kembali
        GDD kelompok Anda, lalu turunkan dua puluh sampai tiga puluh butir
        menurut tiga kelompok pada Bagian~\ref{sec:checklist-fungsional}.
        Isilah bentuk pada Tabel~\ref{tab:checklist-kosong}.

  \item \textbf{Jalankan pemeriksaan otomatis.} Salin peta level Anda ke dalam
        \code{ch16_uji_mandiri.py} dengan bentuk yang sama seperti peta bawaan,
        lalu jalankan:
        \code{python -I ch16_uji_mandiri.py --peta namapeta}. Perbaiki setiap
        tempat yang dilaporkan terjebak atau hanya dapat menurun.

  \item \textbf{Siapkan \istilah{build}.} Bungkus permainan Anda menjadi berkas
        yang dapat dijalankan, lalu ujilah pada komputer lain, bukan pada
        komputer tempat Anda mengembangkannya.

  \item \textbf{Susun lembar \istilah{playtest}.} Isilah bentuk pada
        Tabel~\ref{tab:lembar-playtest}, dan tulislah bagian mana yang paling
        Anda ragukan sebagai tujuan sesi.

  \item \textbf{Laksanakan \istilah{playtest} antar kelompok}
        sekurang-kurangnya dengan dua penguji yang belum pernah melihat karya
        Anda. Ikuti tujuh langkah pada Bagian~\ref{sec:lembar-playtest}, dan
        jangan memberi petunjuk selama sesi berlangsung.

  \item \textbf{Pindahkan temuannya ke \istilah{log bug}} menurut
        Tabel~\ref{tab:log-bug}, dan timbang beratnya menurut
        Tabel~\ref{tab:berat-temuan}. Catatan yang belum dapat diulang jangan
        dipindahkan dahulu; sempitkan dugaan seperti pada
        Bagian~\ref{sec:contoh-bab-16}.

  \item \textbf{Perbaiki temuan berat satu per satu}, dan jalankan kembali
        \istilah{checklist} fungsional sesudah setiap perbaikan. Catat hasil
        setiap putaran pada kolom putaran.

  \item \textbf{Periksa kesesuaian GDD dengan \istilah{build}} menurut
        Tabel~\ref{tab:gdd-build}, dan perbarui GDD untuk setiap perubahan
        rancangan yang disengaja.

  \item \textbf{Susun laporan perbaikan} yang memuat keadaan sebelum dan
        sesudah setiap perbaikan, beserta buktinya.

  \item \textbf{Sajikan demo lima sampai sepuluh menit} menurut susunan pada
        Tabel~\ref{tab:susunan-demo}, dan serahkan \istilah{build},
        \istilah{project source}, lembar \istilah{playtest}, serta
        \istilah{log bug}.

  \item \textbf{Tulis refleksi pemilihan teknologi} dengan menjawab ketiga
        pertanyaan pada Bagian~\ref{sec:presentasi-refleksi}, disertai contoh
        dari pekerjaan kelompok Anda sendiri.
\end{langkah}
\end{praktikumbox}

\begin{table}[htbp]
\centering
\caption{Lembar kerja pengujian prototipe}
\label{tab:lembar-praktikum-16}
{\setlength{\tabcolsep}{4pt}\footnotesize%
\begin{tabularx}{\textwidth}{@{}L{4.0cm} Y@{}}
\toprule
\headrow \textbf{Bagian} & \textbf{Isi} \\
\midrule
Banyak butir \istilah{checklist} & \\[1.3em]
\midrule
Butir yang lulus pada putaran 1 & \\[1.3em]
\midrule
Butir yang lulus pada putaran 2 & \\[1.3em]
\midrule
Temuan berat & \\[1.6em]
\midrule
Temuan sedang dan ringan & \\[1.6em]
\midrule
\addlinespace
\multicolumn{2}{@{}p{\textwidth}@{}}{\footnotesize\itshape
Banyak penguji pada sesi \istilah{playtest}: \dotfill} \\[1.4em]
\addlinespace
\multicolumn{2}{@{}p{\textwidth}@{}}{\footnotesize\itshape
Perbedaan antara GDD dan \istilah{build} yang Anda temukan: \dotfill} \\[1.4em]
\addlinespace
\multicolumn{2}{@{}p{\textwidth}@{}}{\footnotesize\itshape
Perbaikan terpenting yang Anda lakukan, dan buktinya: \dotfill} \\[1.8em]
\bottomrule
\end{tabularx}}
\end{table}

\begin{catatabox}
Langkah kelima adalah langkah yang paling menentukan bagi indikator \ind{a} dan
\ind{b}. Pengujian yang dilakukan dengan saksama akan menghasilkan temuan yang
tidak dapat Anda temukan sendiri, sedangkan pengujian yang disertai petunjuk
akan menghasilkan catatan yang menyenangkan tetapi tidak berguna.
\end{catatabox}
